% ==========================================================================
%  Chapter 9 — The v2 Redesign Project
% ==========================================================================
\chapter{The v2 Redesign Project}
\label{ch:v2redesign}

\epigraph{%
  A good architecture is one that allows the system to be easily changed.%
}{Robert C. Martin}

\section{Project Overview}
\label{sec:v2-overview}

The v2 redesign of User-Mode Linux is a ground-up rearchitecture of
the \UML{} backend layer, validation infrastructure, and operational
tooling.  The work lives in the \texttt{mjbommar/linux} kernel tree
(branch \texttt{next}), based on \texttt{torvalds/master} at
v7.1-rc5.  It was developed between late 2025 and May 2026 by
Michael~J.~Bommarito~II with substantial AI-assisted tooling
(Anthropic Claude).

The tree is organized around four branches:

\begin{description}[style=nextline,leftmargin=2em]
  \item[\texttt{master}]
    Tracks \texttt{torvalds/master}; the baseline.
  \item[\texttt{next}]
    The integration branch.  All seven series-level commits live here,
    rebased onto \texttt{master} at each Linus pull.
  \item[\texttt{fixes}]
    Hotfixes that will be cherry-picked into \texttt{next} and submitted
    independently.
  \item[\texttt{umlctl-deploy}]
    Deployment-tested snapshots of \texttt{tools/uml/uml-launcher/}
    for operators who need a stable \texttt{umlctl} without tracking
    the kernel tree.
\end{description}

\noindent
The development history comprises roughly 1,236 development commits,
compressed during rebase into seven clean series-level commits on
\texttt{next}.  Each series is independently reviewable and
bisectable.  The total diff against \texttt{torvalds/master} spans
approximately 248,000 insertions across 1,005 files, of which roughly
134,500 lines are documentation, 54,000 are tooling and selftests,
33,400 are the \KVMvTwo{} backend, and 25,900 are the backend-ops
abstraction and core infrastructure.


\section{Design Goals}
\label{sec:v2-goals}

The redesign is driven by seven concrete, measurable goals drawn from
the project vision document:

\begin{enumerate}[leftmargin=2em]
  \item \textbf{Near-native syscall overhead on the \KVM{} backend.}
    The target is approximately \SI{100}{\nano\second} per guest syscall
    when the host exposes \texttt{/dev/kvm}.  The achieved result is
    approximately 90~cycles ($\sim$\SI{25}{\nano\second} on modern
    hardware) for LSTAR-gadget-covered syscalls---exceeding the
    original goal by roughly $4\times$.

  \item \textbf{Full kernel instrumentation.}
    KASAN, KMSAN, KCSAN, KFENCE, KCOV, kprobes, ftrace, and BPF JIT
    must all be composable from a single tree.  Instrumentation is
    controlled by Kconfig profiles and static keys, not by separate
    builds or out-of-tree patches.

  \item \textbf{Deterministic time-travel and record/replay.}
    The record/replay infrastructure captures non-deterministic events
    (SYSCALL exits, RDTSC results, SIGALRM deliveries) during
    execution and replays them faithfully.  The time-travel profile
    provides deterministic clock semantics.

  \item \textbf{Snapshot/forkserver startup under \SI{50}{\milli\second}.}
    The achieved result is \SI{38.9}{\micro\second} capture and
    \SI{13}{\micro\second} median restore---roughly $1{,}300\times$
    and $75\times$ faster than the original targets, respectively.

  \item \textbf{Three explicit profiles with inverted defaults.}
    Research, fuzz, and sandbox profiles ship with different backend
    selections, different instrumentation, and different host attack
    surface.  The profile system spans nine named configurations
    (Table~\ref{tab:v2-profiles}), each built from the same tree
    with a single \texttt{make ARCH=um O=\ldots} invocation.

  \item \textbf{Runtime-flippable hooks.}
    Static keys (\texttt{DECLARE\_STATIC\_KEY\_FALSE}) gate
    instrumentation hooks at the binary level.  A
    \texttt{prod-with-hooks} build compiles all hook sites but
    patches them to NOPs at boot; a debugfs write or a crash handler
    can enable them without rebuilding.

  \item \textbf{Backend abstraction with swappable trap mechanisms.}
    \texttt{struct um\_backend\_ops} defines a contract for backend
    lifecycle, guest process creation, syscall interception, context
    switching, and memory mapping.  The ptrace, seccomp, and \KVMvTwo{}
    backends implement this interface.  New backends can be added
    without modifying core \UML{} code.
\end{enumerate}

\begin{table}[ht]
\centering
\caption{Profile matrix: nine named configurations from one tree.}
\label{tab:v2-profiles}
\small
\begin{tabularx}{\linewidth}{@{}p{0.16\linewidth}p{0.16\linewidth}p{0.20\linewidth}p{0.14\linewidth}Y@{}}
\toprule
Profile & Backend & Hooks & Sanitizers & Purpose \\
\midrule
\texttt{prod-fast}       & KVM $>$ seccomp & none            & none          & Minimum overhead \\
\texttt{prod-with-hooks} & KVM $>$ seccomp & compiled, off   & none          & Post-failure observability \\
\texttt{research}        & seccomp         & trace, kprobes  & KASAN, UBSAN  & Debug and exploit reproduction \\
\texttt{fuzz}            & seccomp         & KCOV, snapshot  & KASAN         & syzkaller integration \\
\texttt{fuzz-deep}       & seccomp         & KCOV, snap, replay & KASAN, KCSAN & Race and replay investigations \\
\texttt{sandbox}         & seccomp-only    & none compiled   & none          & Minimized TCB \\
\texttt{library}         & direct call     & n/a             & opt.\ KASAN   & LKL-style harness \\
\texttt{embedded}        & ptrace          & none            & none          & Old-host fallback \\
\texttt{time-travel}     & seccomp         & trace, TT       & KASAN         & Deterministic clock / replay \\
\bottomrule
\end{tabularx}
\end{table}


\section{What We Are Not Building}
\label{sec:v2-nongoals}

The design boundary is explicit.  The v2 redesign is not:

\begin{description}[style=nextline,leftmargin=2em]
  \item[A kernel reimplementation (gVisor).]
    gVisor's Sentry~\cite{gvisor} reimplements the Linux syscall
    interface in Go.  \UML{} \emph{is} the Linux kernel.  Every
    scheduler decision, every lock acquisition, every VFS path follows
    the same code as a kernel running on x86 hardware.

  \item[A unikernel (Unikraft, MirageOS).]
    Unikernels~\cite{kuenzer2021,madhavapeddy2013} compile application
    and kernel into a single-address-space image, sacrificing process
    isolation for boot speed.  \UML{} maintains genuine user/kernel
    separation.

  \item[A microVM (Firecracker, Cloud Hypervisor).]
    Firecracker~\cite{agache2020} and Cloud Hypervisor~\cite{cloud-hypervisor}
    are minimal VMMs optimized for serverless cold-start.  \UML{}'s
    value proposition is kernel research and debugging fidelity, not
    cloud packaging.

  \item[A container runtime.]
    Containers share the host kernel.  \UML{} runs its own kernel,
    with its own scheduler, its own VFS, and its own memory management.
    A bug in a container's workload can exploit the host kernel; a bug
    in a \UML{} guest's workload can only exploit the \UML{} kernel.

  \item[A hypervisor.]
    \UML{} is a guest, not a host.  It does not manage hardware, does
    not provide a VMM API, and does not schedule other virtual machines.
    When \UML{} uses \KVM{}, it is a \KVM{} \emph{consumer}, not a
    \KVM{} replacement.
\end{description}

\noindent
The clearest way to state the design intent: \UML{} is the full Linux
kernel, running in userspace, with pluggable trap mechanisms and
composable instrumentation.  Nothing more, nothing less.


\section{The Series Breakdown}
\label{sec:v2-series}

The seven series-level commits are organized by subsystem scope,
review audience, and dependency order.  Table~\ref{tab:v2-series}
summarizes the series; the subsections below provide detail.

\begin{table}[ht]
\centering
\caption{Seven series-level commits on \texttt{torvalds/master}.}
\label{tab:v2-series}
\small
\begin{tabularx}{\linewidth}{@{}r p{0.30\linewidth} r r Y@{}}
\toprule
\# & Series & Files & Insertions & Subsystem \\
\midrule
1 & BPF JIT stub + x86 hygiene         &   2 &     123 & BPF \\
2 & KMSAN early-shadow callback         &   2 &      43 & mm/kmsan \\
3 & notrace on kthread/smpboot          &   2 &      31 & tracing/ftrace \\
4 & Backend ops abstraction (RFC)       & 165 &  25,869 & arch/um \\
5 & \KVMvTwo{} backend                  &  52 &  33,374 & arch/um + KVM \\
6 & \texttt{umlctl} + kselftests        & 364 &  54,000 & tools/ \\
7 & Documentation                       & 420 & 134,539 & Documentation/ \\
\midrule
  & \textbf{Total}                      & \textbf{1,005} & \textbf{247,979} & \\
\bottomrule
\end{tabularx}
\end{table}


\subsection{Series 1: BPF JIT Stub and x86 Hygiene}
\label{sec:v2-series1}

\textbf{Commit:} \texttt{c1084ea01e06} ---
\texttt{um: bpf: add JIT stub and x86 hygiene for UML}

Two files changed, 123 insertions.  This series adds
\texttt{arch/um} BPF JIT stubs and fixes x86 BPF JIT compilation
under \UML{}, where direct code generation targets the host's JIT
engine.  The change is small, self-contained, and reviewable by the
BPF maintainers independently of any \UML{} redesign.


\subsection{Series 2: KMSAN Early-Shadow Callback}
\label{sec:v2-series2}

\textbf{Commit:} \texttt{be9dbd18e75e} ---
\texttt{mm/kmsan: add arch early-shadow callback for UML physmem layout}

Two files changed, 43 insertions.  \UML{}'s physical memory
(\texttt{physmem}) is a single contiguous \syscall{mmap} region, not
the standard \texttt{pfn-to-page} layout that KMSAN's generic
initialization expects.  This series adds an architecture callback in
\texttt{kmsan\_init\_shadow()} so \UML{} can provide its
shadow/origin mappings from its own address space.  The change touches
\texttt{mm/kmsan/} and is reviewable by the KMSAN maintainer
independently.


\subsection{Series 3: notrace on Signal-Adjacent Thread Functions}
\label{sec:v2-series3}

\textbf{Commit:} \texttt{8301623cf952} ---
\texttt{kthread/smpboot: add notrace to signal-adjacent thread functions}

Two files changed, 31 insertions.  Under
\texttt{-fpatchable-function-entry} (used by ftrace on architectures
including \UML{}), \texttt{kthread()} and \texttt{smpboot\_thread\_fn()}
are instrumented, but they run in contexts where the ftrace trampoline
may not be safe (signal delivery, early boot).  The \texttt{notrace}
annotation is a generic improvement that benefits any architecture
with patchable function entry, not only \UML{}.


\subsection{Series 4: Backend Ops Abstraction (Tentpole RFC)}
\label{sec:v2-series4}

\textbf{Commit:} \texttt{86b74b8993d0} ---
\texttt{um: backend ops abstraction + SMP + TLB sync + infrastructure}

This is the tentpole series: 165 files changed, 25,869 insertions,
approximately 932 deletions.  It introduces the core architectural
change that the rest of the redesign depends on.

\paragraph{struct um\_backend\_ops.}
A typed dispatch table with operations for backend lifecycle (init,
cleanup), guest process management (create, fork, exec, exit),
syscall interception (trap, return), context switching, and memory
mapping.  The existing ptrace and seccomp code paths are refactored
through this table.

\paragraph{SMP support.}
Virtual CPUs are mapped to host threads with per-CPU dispatch.
Inter-processor interrupt (IPI) emulation uses \texttt{eventfd} on
the host.  TLB synchronization is rewritten with a deferred-free
mechanism, \texttt{sync\_tlb\_range} pattern, and
\texttt{page\_table\_lock} integration.

\paragraph{PTE rework.}
PTE bits are realigned to x86 hardware paging format, enabling a
single-store \texttt{set\_pte} that eliminates the TDP walker race
present in the old multi-store sequence.  The \texttt{pte\_read}/\texttt{pte\_exec}
fix (adopted upstream as \texttt{cd4126d48f7f}) is included.

\paragraph{Infrastructure.}
Additional components include hostfs io\_uring writeback substrate,
section-split infrastructure for \texttt{.text.frozen},
snapshot/forkserver worker lifecycle management, deferred TLB batch
for \texttt{mm/mmu\_gather.c}, per-CPU RSS drift correction for
\texttt{mm/mmap.c}, and a CFS worker-detach helper in
\texttt{kernel/sched/core.c}.


\subsection{Series 5: KVM v2 Backend}
\label{sec:v2-series5}

\textbf{Commit:} \texttt{12b19d012f9f} ---
\texttt{um: kvm-v2: add KVM-based backend for User-Mode Linux}

52 files changed, 33,374 insertions.  This series adds the complete
\KVM{}-based execution backend.  The guest runs inside a \KVM{} VM
with EPT/VT-x, replacing the ptrace/seccomp signal-mediated syscall
path with \texttt{VMRESUME}/\texttt{VMEXIT}.

Key components of \umlpath{arch/um/backend/kvm-v2/}:

\begin{description}[style=nextline,leftmargin=2em]
  \item[\texttt{vcpu.c}] Per-vCPU lifecycle, \texttt{KVM\_RUN}
    dispatch loop, SREGS synchronization.
  \item[\texttt{syscall\_trap.c}] IDT exception dispatch covering all
    x86 vectors (\#DE, \#DB, \#BP, \#OF, \#UD, \#NM, \#DF, \#TS,
    \#NP, \#SS, \#GP, \#PF, \#MF, \#AC, \#XM).
  \item[\texttt{exception.c}] IDT/GDT/TSS/handler page setup.
  \item[\texttt{context.c}] VM lifecycle and physmem memslot management.
  \item[\texttt{lstar\_gadget.S}] In-guest \texttt{SYSCALL} trampoline
    that handles fast-path syscalls without a \texttt{VMEXIT},
    achieving $\sim$90-cycle latency.
  \item[\texttt{snapshot.c}] KVM-aware snapshot with capture/restore
    for regs, sregs, and memslot state.
  \item[\texttt{record.c}] Record/replay state machine with static-key
    hot-path gating.
  \item[\texttt{test\_snapshot.c}, \texttt{test\_record.c}] KUnit test
    suites (3/3 and 7/7 cases, respectively).
\end{description}

\noindent
The series also preserves the v1 \KVM{} backend as
\umlpath{arch/um/backend/kvm-v1-archive/} for reference, and includes
contract tests in \umlpath{arch/um/backend/contract/}.


\subsection{Series 6: umlctl and Kselftests}
\label{sec:v2-series6}

\textbf{Commit:} \texttt{dba578471251} ---
\texttt{um: add umlctl launcher tool + kselftests}

364 files changed, 54,000 insertions.  This series adds two major
components.

\paragraph{umlctl (\umlpath{tools/uml/uml-launcher/}).}
A Rust CLI for managing \UML{} instances: create, start, stop,
\texttt{ps}, logs, metrics, and \texttt{up} (declarative deployment
from TOML Umlfiles).  The crate comprises 65 Rust source files
totaling approximately 22,000 lines, plus profiles, examples, and
scripts.  Key subcommands include:

\begin{itemize}[leftmargin=2em]
  \item \texttt{umlctl mission} --- Six-phase acceptance gate with
    JSON scoreboard and binary verdict ($\sim$15 minutes wall-clock).
  \item \texttt{umlctl up} --- Declarative instance launch from TOML.
  \item \texttt{umlctl build} --- Kernel build with profile selection.
  \item Cgroup v2 per-instance subgroup management and preflight
    resource verification.
\end{itemize}

\paragraph{kselftests (\umlpath{tools/testing/selftests/um/}).}
53 test suites spanning 316 files.  Categories include:

\begin{itemize}[leftmargin=2em]
  \item \textbf{Substrate parity:} \texttt{cpython-parity},
    \texttt{cpython-tier0}, \texttt{cpython-full}.
  \item \textbf{Backend tests:} \texttt{kvm-smoke},
    \texttt{kvm-isolation}, \texttt{kvm-bounds}, \texttt{kvm-mm-smoke}.
  \item \textbf{Snapshot/replay:} \texttt{snapshot-smoke},
    \texttt{snapshot-kvm-smoke}, \texttt{kvm-record-smoke},
    \texttt{kvm-snapshot-bench}, \texttt{kvm-record-clock-bench}.
  \item \textbf{SMP stress:} \texttt{mt-mmap-stress},
    \texttt{fork-tree-3level}, \texttt{template-pause-*}.
  \item \textbf{Performance:} \texttt{bench}, \texttt{perf-getpid},
    \texttt{perf-py-startup}, \texttt{pool-bench}, \texttt{net-bench}.
  \item \textbf{Instrumentation:} \texttt{ftrace-smoke},
    \texttt{kmsan-smoke}, \texttt{kprobes-stress}, \texttt{hooks-flip}.
  \item \textbf{Soak:} 11 TOML-templated workloads for sustained
    validation (memcheck, iocheck, stress-ng, tier1-pylibs,
    django-loopback, and more).
\end{itemize}


\subsection{Series 7: Documentation}
\label{sec:v2-series7}

\textbf{Commit:} \texttt{f8d5d3df25a8} ---
\texttt{Documentation/virt/uml: add comprehensive UML documentation}

420 files changed, 134,539 insertions.  This series adds operator
documentation for all \UML{} subsystems (backends, kprobes, ftrace,
KMSAN, debugfs, snapshot/ELF-core, section-split, profiles), the
complete redesign planning tree (vision, architecture, workstreams,
sequencing, decisions log, risk register), the upstream submission
queue with per-series submission notes, and CI workflows for
checkpatch, \KVM{} probe, uml-launcher build, and profile boot
validation.


\section{Validation Methodology}
\label{sec:v2-validation}

The v2 redesign is validated through a multi-tier methodology that
progresses from unit tests through full-workload soak runs.  The
design principle is that each tier must pass before the next is
attempted, and the entire ladder is automated through
\texttt{umlctl mission}.

\subsection{The Validation Ladder}

\begin{enumerate}[leftmargin=2em]
  \item \textbf{Tooling builds.}  The Rust \texttt{umlctl} crate and
    all selftests compile cleanly.

  \item \textbf{Kernel builds.}  Every named profile compiles and
    links under \texttt{make ARCH=um O=\ldots}.

  \item \textbf{Internal contracts.}  KUnit suites pass: snapshot 3/3,
    record/replay 7/7, marshal/byte-shape checks.

  \item \textbf{Substrate parity gate.}  The CPython test suite runs
    on both backends with identical results: PASS=25, FAIL=3, XFAIL=3.
    The 21 modules in the cpython-parity gate pass identically on
    \KVMvTwo{} and seccomp, confirming functional equivalence.

  \item \textbf{Process and syscall classes.}  Dedicated selftests
    exercise fork, exec, mmap, signals, dynamic loader, and
    \texttt{rt\_sigreturn} isolation.

  \item \textbf{Performance gates.}  Micro (getpid-loop),
    application (Python startup), and stress (mt-mini SMP T=8 ncpus=4)
    tiers, plus snapshot benchmark (capture \SI{38.9}{\micro\second},
    restore \SI{13}{\micro\second}).

  \item \textbf{Real workload boot.}  Python and stress workloads
    complete successfully on both backends.

  \item \textbf{\texttt{umlctl mission} acceptance gate.}
    Six phases, $\sim$15 minutes wall-clock, 100/100 soak iterations,
    binary PASS/FAIL verdict with JSON scoreboard.

  \item \textbf{24-hour long-tail soak.}  The remaining wall-clock
    evidence for milestone candidates.
\end{enumerate}

\subsection{The \texttt{umlctl mission} Gate}

The \texttt{umlctl mission} command is the canonical acceptance
artifact.  A single invocation runs six fail-fast phases:

\begin{enumerate}[leftmargin=2em]
  \item \textbf{KUnit} ($\sim$5\,s) --- 10/10 cases PASS (snapshot 3 +
    record 7).
  \item \textbf{Bench} ($\sim$3\,s) --- Capture \SI{35}{\micro\second},
    median restore \SI{13}{\micro\second}.
  \item \textbf{Substrate} ($\sim$4\,s) --- cpython-tier0 PASS.
  \item \textbf{Host resources} ($\sim$1\,s) --- Environment knobs applied.
  \item \textbf{Diverse soak} ($\sim$14\,min) --- 100/100 PASS across
    memcheck, iocheck, stress-ng, tier1-pylibs, and
    django-loopback; zero panics, zero SIGBUS, zero high-cr2.
  \item \textbf{Diagnostic} ($<$1\,s) --- Kernel SHA, host metadata,
    hugepage pool sizing, cgroup v2 state captured.
\end{enumerate}

\noindent
Any single regression in any phase produces
\texttt{MISSION\_FAILED phase=\textit{N}} on stderr and a non-zero
exit.  The gate explicitly fails on panics, ``Kernel mode signal~7''
(\texttt{/dev/shm} exhaustion), and ``high-cr2'' (the residual
user-fault class).

\subsection{Third-Party and Real-World Tests}

Beyond the in-tree selftests, the validation methodology includes:

\begin{itemize}[leftmargin=2em]
  \item \textbf{Tier 1 Python libraries:} NumPy, requests, Flask,
    SQLAlchemy, and other commonly-used packages exercised during soak.
  \item \textbf{Django loopback:} A Django test server with HTTP
    request/response validation.
  \item \textbf{LTP (Linux Test Project):} Curated subset with
    \UML{}-appropriate runner refinements (in progress).
  \item \textbf{stress-ng:} CPU, memory, and I/O stress patterns.
\end{itemize}


\section{The CPython Bug Discovery}
\label{sec:v2-cpython-bug}

During substrate parity validation, \UML{} uncovered a latent bug in
CPython 3.14's \texttt{Py\_FinalizeEx} teardown sequence.  The bug
manifested as a segmentation fault during interpreter shutdown, caused
by a race condition in \texttt{\_Py\_Dealloc} where the thread state
(\texttt{tstate}) pointer is NULL-dereferenced after
\texttt{PyStdPrinter\_Type}'s deallocation hook runs.

On hardware and under conventional hypervisors, the race is nearly
invisible: the timing window is too narrow to hit reliably.  \UML{}'s
slower execution---particularly under the seccomp backend---widens
the window enough that the race triggers reproducibly.  The \KVMvTwo{}
backend, being faster, hits it less frequently but still observably.

\paragraph{The fix.}
Commit \texttt{40203d69dff1} adds a five-condition intercept in
\umlpath{arch/um/kernel/trap.c::segv()} that detects the specific
\texttt{tstate}-NULL dereference pattern during \texttt{Py\_FinalizeEx}
and converts it to a clean \texttt{exit(1)} instead of a kernel
panic.  The fix was validated with zero SEGV occurrences per run on
both backends.

\paragraph{The broader lesson.}
This discovery demonstrates a category of value that no other
virtualization technology provides.  Hardware runs too fast to expose
the race.  Containers share the host kernel and cannot alter execution
timing.  gVisor reimplements the syscall layer and may not reproduce
the same ordering.  Only \UML{}---the real kernel, running at
controllable speed, with full debugging---can simultaneously reproduce
the bug, diagnose the root cause, and validate the fix.  It is
precisely the kernel-fidelity-plus-debuggability combination that
justifies \UML{}'s continued existence.


\section{Upstream Submission Plan}
\label{sec:v2-upstream}

The upstream strategy is deliberately staged to minimize review burden
and maximize the probability of acceptance.

\subsection{Phase 1: Independent Subsystem Patches (Series 1--3)}

Series 1 (BPF JIT stub), Series 2 (KMSAN early-shadow), and
Series 3 (notrace on kthread/smpboot) are small, self-contained
patches that touch code outside \umlpath{arch/um/}.  They are
reviewable by the respective subsystem maintainers (BPF, mm/kmsan,
tracing/ftrace) without requiring any understanding of the \UML{}
redesign.  These are submitted first to establish a track record and
to get the cross-subsystem prerequisites in place.

\subsection{Phase 2: The Tentpole RFC (Series 4)}

Series 4, the backend-ops abstraction, is the architectural
centerpiece.  It refactors existing code through a new interface
without changing observable behavior for existing backends.  The RFC
submission is accompanied by a detailed cover letter explaining:

\begin{itemize}[leftmargin=2em]
  \item Why the abstraction is needed (multiple backends already exist;
    making the contract explicit is a cleanup).
  \item What the interface looks like (\texttt{struct um\_backend\_ops}).
  \item How existing backends are migrated (ptrace and seccomp
    refactored through the ops table).
  \item What the abstraction enables (the \KVMvTwo{} backend, future
    backends).
\end{itemize}

\noindent
This is the series that must land before anything else can follow.

\subsection{Phase 3: KVM Backend and Tooling (Series 5--7)}

Series 5 (the \KVMvTwo{} backend), Series 6 (\texttt{umlctl} and
kselftests), and Series 7 (documentation) depend on Series 4.  They
are submitted after the backend-ops abstraction is accepted or has
received sufficient positive review signal.  The cover letter for
Series 5 includes the complete benchmark data
(Appendix~\ref{app:benchmarks}), the \texttt{umlctl mission}
acceptance gate results, and instructions for reviewers to reproduce
the validation locally.

\subsection{The LKML Strategy}

The submission targets the \texttt{linux-um} mailing list with
CC to relevant subsystem maintainers.  Key considerations:

\begin{itemize}[leftmargin=2em]
  \item AI-assisted development is disclosed per the kernel's
    \texttt{coding-assistants.html} policy~\cite{coding-assistants}.
  \item Each series carries a \texttt{Signed-off-by} from the human
    author and a \texttt{Co-authored-by} from the AI tool.
  \item The submission queue document
    (\umlpath{upstream-patches/SUBMISSION-QUEUE.md}) tracks the
    status of each series.
\end{itemize}


\section{Lessons Learned}
\label{sec:v2-lessons}

Several lessons emerged from the development process that may be
useful to future large-scale kernel rearchitecture efforts.

\paragraph{1. The backend contract must be explicit and tested.}
The single most valuable architectural decision was defining
\texttt{struct um\_backend\_ops} early and writing contract tests
(\umlpath{arch/um/backend/contract/}) that verify each backend
against the same behavioral specification.  Every subsequent bug
was diagnosed faster because the contract narrowed the search space.

\paragraph{2. Substrate parity is the right acceptance criterion.}
Using CPython's test suite as the primary acceptance criterion---rather
than microbenchmarks or unit tests alone---caught bugs that no
synthetic test would have found.  The CPython teardown race
(Section~\ref{sec:v2-cpython-bug}) is the most dramatic example, but
subtler issues in signal delivery, mmap semantics, and dynamic loader
behavior were also caught by substrate parity testing.

\paragraph{3. State audits prevent regression spirals.}
The v1 \KVM{} backend was retired because of accumulated state
ownership confusion between the \UML{} kernel and the \KVM{} host.
The v2 design was informed by a formal state audit
(\umlpath{02-workstreams/D-kvm-backend/state-audit/}) that mapped
every piece of shared state, its owner, its synchronization mechanism,
and its failure mode.  This audit prevented the same class of bugs
from recurring.

\paragraph{4. The investigation discipline pays for itself.}
The project memory records a discipline: ``run the cheapest
dispositive control FIRST; never declare `architectural' or ship
userspace workarounds without it.''  Every time this discipline was
followed, it saved hours.  Every time it was violated (the pool
physmem isolation episode, the SIGALRM/mmap-FIXED deadlock), it cost
days.

\paragraph{5. Compress early, compress often.}
The 1,236 development commits exist because iterative development
requires iterative commits.  But for upstream submission, only the
logical decomposition matters.  The rebase from 1,236 commits to 7
series was done once, at the end, when the design was stable.
Attempting to maintain a clean commit history during active
development would have been a net drag on velocity.
